\section{Answer and evidence boundary}

There are useful transfers beyond familiar consequences of RH and Dirichlet GRH. The strongest conclusion supported by this investigation has two parts. First, a previous ProveIt report contains a complete extra arithmetic argument for a classification that was still conditional on an unsearched interval in the cited literature. Its decisive certificates reproduce successfully. Second, the internal proof machinery has genuinely reusable pieces, but the passage to a new nonabelian zero theorem still requires an arithmetic comparison not supplied by the repository or by the incoming reports.

\subsection{A result that can be endorsed now}

The exponent-six completion in \emph{grh-exponent-six-and-transfers.zip} proves that exactly 432 imaginary quadratic fields have ideal class group of exponent six. Aiken--Ross--Shvydkoy--Xue's Theorem 1.3(5) still leaves a fundamental-discriminant interval to exclude even under ERH. The previous report eliminates that interval by an explicit-formula test, a collision between two integral norm equations, and an exact residue sieve. Here its mathematical reduction was reviewed, its interval-arithmetic test rerun, and both independent complete divisor searches reproduced. This is a validation and re-exposition of earlier work, not a claim of discovering the classification during the present turn.\cite{priorSix,EKN,ARSX}

The extension from fields to orders then gives 1,236 orders of exact exponent six. The field list and the order-table identity were checked. The underlying published exhaustive enumeration through $3.1\cdot10^{20}$ was not rerun; it is used as a published computational theorem. Section~\ref{sec:census} makes all dependencies explicit.

This application uses the exported quadratic-GRH theorem together with a substantial additional argument. It differs from importing a ready-made conditional consequence, and it also differs from directly transplanting the source's unitary theta-lift construction. Both distinctions matter to the requested scope.

\subsection{New work in this synthesis}

The following additions were not found in the 26 inspected archives. Their mathematical ingredients have classical antecedents, so no claim of worldwide priority is made.
\begin{enumerate}
\item An explicit trace frame over the real $S_3$ splitting field of $T^3-4T+1$, together with a cyclotomic tensor construction and a precise local Frobenius compatibility. This realizes one local input for a nonabelian extension; it is more concrete than merely specifying a desired prime statistic.
\item A joint positive-averaging criterion: several individually imperfect averages detect every pole if their Mellin multipliers have no common zero. An integer-coefficient two-filter example comes with an exact polynomial identity and a sharp obstruction when a common zero exists.
\item A positive-ramp $L^p$-to-pointwise estimate with optimal exponent, including a countermodel with uniformly bounded positive real coefficients supported on the integers. It explains why direct averaged Landau arguments are valuable: one should not insist on recovering the full pointwise error first.
\end{enumerate}

\subsection{What the source proves}

The public endpoints in \src{Statements}{Statements.lean} and \src{Main}{Main.lean} concern the Riemann zeta function and Dirichlet characters over $\Q$. In particular, the theorem \texttt{LFunction\_eq\_zero\_re} states that a Dirichlet $L$-function zero in $0<\Re s<1$ has real part $1/2$. It does not quantify over all Hecke characters of all number fields, all Artin representations, or all automorphic representations.\cite{source}

Independent public checks report successful compilation and axiom checks of both shortened editions; the compact project also has a reported complete Lean kernel replay. These external reports support taking the supplied theorem as the starting point. The source examination in this session is not an independent rerun of those complete checks.\cite{Sutherland,Rubinstein}

\subsection{A baseline that avoids repeated reductions}

The earlier reports already develop precise nonabelian detectors for cubic and quintic fields, $S_n$, $A_5$, $A_6$, Sylow subgroups, cyclic generators, reflection groups, and $\operatorname{PSL}_2$. They also contain reciprocal convexity, fourth-power sampling, positive-part Mellin criteria, moving-box averages, torsion selection, and ramification-aware Selmer warnings. Appendix~A and the accompanying claim matrix record all 26 packages.\cite{priorCollection}

The fourth-power sampling idea is therefore explicitly credited to the prior reports and their cited antecedent, rather than counted as new. Likewise the positive-part Mellin lemma below is reused from the previous mean-transfer report. The contribution of the present analytic sections is the joint-filter repair and the sharp positive-ramp specialization.\cite{priorSampling,priorMean,Setiawan}
